\lecture{7}
\title{DT Fourier Transform}

\begin{document}

\subsection{Lecture: The DT Fourier Transform}

Let's do it all again for DT.
You might take the following as a prototypical example.
\[ f[n] := \begin{cases}
1 & \text{ if } n \in \{-2, -1, 0, 1, 2\}. \\
0 & \text{ otherwise.}
\end{cases} \]
The goal, as before, is to express $f[n]$ as a sum (or rather, integral) of complex exponentials,
even though $f[n]$ is aperiodic.

We'll just cut straight to the punchline this time.

\begin{theorem}{DT Fourier Transform}
    Given a DT signal $f[n]$, define its Fourier transform to be
    $F(\Omega) := \sum_{n = -\infty}^{\infty} f[n] e^{-j \Omega n}$.
    It then follows that
    $f[n] = \frac{1}{2\pi} \int_{0}^{2\pi} F(\Omega) e^{j \Omega n} \, \mathrm{d}\Omega$.
\end{theorem}

\begin{remark}
    Even though the signal $f[n]$ is discrete and possibly aperiodic,
    the Fourier transform $F(\Omega)$ is continuous and guaranteed periodic
    with period $2\pi$.
\end{remark}

\begin{proof}
    Our strategy is the same as last time:
    consider signals $f_N[n]$ that are extensions of $f[n]$ to have period $N$,
    take the DT Fourier series of $f_N[n]$, and investigate the series for $N \to \infty$.

    The DT Fourier series for $f_N[n]$ looks like:
    \[ f_N[n] = \sum_{k \in \langle N \rangle} F_N[k] e^{jn \textcolor{purple}{(2\pi k / N)}}
    ~ \text{ where } ~
    F_N[k] = \frac{1}{N} \sum_{n \in \langle N \rangle} f_N[n] e^{-jn \textcolor{purple}{(2\pi k / N)}}. \]
    Here, $\langle N \rangle$ denotes a set of $N$ consecutive integers roughly centered at zero.
    Remember that for DT Fourier series, $F_N[k]$ is a finite array of length $N$,
    or equivalently, an infinite array with period $N$.

    Let's follow the recipe from last time: take $F(\Omega) := \lim_{N \to \infty} N \cdot F_N[k]$,
    where the \emph{continuous} frequency variable $\textcolor{purple}{\Omega := 2\pi k / N}$
    is held fixed as $N \to \infty$ (so the index $k$ grows in proportion to $N$).  Then:
    \[ F(\Omega) = \lim_{N \to \infty} N \cdot \left(
        \frac{1}{N} \sum_{n \in \langle N \rangle} f_N[n] e^{-jn \textcolor{purple}{(2\pi k / N)}}
    \right)
    = \lim_{N \to \infty} \sum_{n \in \langle N \rangle} f_N[n] e^{-jn \textcolor{purple}{\Omega}}
    = \sum_{n = -\infty}^{\infty} f[n] e^{-jn \textcolor{purple}{\Omega}}. \]
    Just as promised.
    Now make the parametrization $\{\Omega_k\}_{k \in \langle N \rangle}$ where
    $\textcolor{purple}{\Omega_k := 2\pi k / N}$.
    Noting that $\Delta \Omega := \Omega_{k + 1} - \Omega_k = 2\pi / N$
    and sending $N \to \infty$, the Fourier series becomes:
    \begin{align*} f[n] = \lim_{N \to \infty} f_N[n]
    & = \lim_{N \to \infty} \sum_{k \in \langle N \rangle}
    F_N[k] e^{jn \textcolor{purple}{(2\pi k / N)}} \\
    & = \lim_{N \to \infty} \frac{1}{2\pi} \sum_{k \in \langle N \rangle}
    F(\Omega_k) e^{jn\textcolor{purple}{\Omega_k}} \Delta \Omega
    = \frac{1}{2\pi} \int_{\langle 2\pi \rangle}
    F(\Omega) e^{jn \textcolor{purple}{\Omega}} \, \mathrm{d}\Omega. \end{align*}
    Just as promised.  Here, $\langle 2\pi \rangle$ refers to
    any interval of length $2\pi$; it appears because the set of frequencies
    $\{\Omega_k\}_{k \in \langle N \rangle}$ covers an interval of length $2\pi$
    in the limit $N \to \infty$.  ~ $\blacksquare$
\end{proof}

\begin{remark}
    Given a DT signal $f[n]$, consider the ``staircase'' CT signal $g(t) := f[\lfloor t \rfloor]$.
    Then the following claim (perhaps surprisingly!) is false.
    \begin{quote}
        \textbf{False Claim.} The CT Fourier transform $G(\omega)$ of $g(t)$ is the same as
        the DT Fourier transform $F(\Omega)$ of $f[n]$.
    \end{quote}
    To see why this claim fails, compute $G(\omega)$ one staircase step at a time:
    \[ G(\omega) = \sum_{n = -\infty}^{\infty} f[n] \int_{n}^{n + 1} e^{-j \omega t} \, \mathrm{d}t
    = \sum_{n = -\infty}^{\infty} f[n] e^{-j \omega n} \left[ \frac{1 - e^{-j\omega}}{j\omega} \right]
    = F(\Omega) \left[ \frac{1 - e^{-j\omega}}{j\omega} \right].
    ~~~~~ (\Omega :=: \omega) \]
    So the CTFT $G(\omega)$ is the DTFT $F(\Omega)$ \emph{times} an extra decay factor of
    $\left[\frac{1 - e^{-j\omega}}{j\omega}\right]$.

    This factor is also why the synthesis equations differ:
    the CTFT $G(\omega)$ decays, so it's aperiodic and thus synthesized by $\int_{-\infty}^{\infty}$.
    In contrast, $F(\Omega)$ is $2\pi$-periodic and thus synthesized by $\int_0^{2\pi}$.
\end{remark}

\begin{problem}
    Compute the DT Fourier transform of the following DT signal, where $S$ is a positive integer.
    \[ f[n] := \begin{cases}
    1 & \text{ if } |n| \leq S. \\
    0 & \text{ otherwise.}
    \end{cases}\]
\end{problem}

\begin{solution}
    The analysis equation gives us:
    \[ F(\Omega) = \sum_{n = -S}^{S} \left(e^{-j\Omega}\right)^n
    = e^{jS\Omega} \left(\frac{1 - e^{- (2S + 1)j\Omega}}{1 - e^{-j \Omega}}\right)
    = \frac{
        e^{(S + \frac{1}{2})j \Omega} - e^{-(S + \frac{1}{2})j \Omega}
    }{
        e^{\frac{1}{2}j \Omega} - e^{-\frac{1}{2}j \Omega}
    } = \frac{\sin\left(\left(S + \frac{1}{2}\right)\Omega\right)}{\sin\left(\frac{1}{2}\Omega\right)}. \]
    Here, we simplify a finite geometric series,
    then multiply the numerator and denominator by $e^{j \Omega / 2}$
    to reach a quotient of sines. ~ $\blacksquare$
\end{solution}

\subsection{Lecture and Recitation: Properties of the DT Fourier Transform}

Here's what the Fourier transforms of the $f[n]$ signal from the previous example
look like for different values of $S$.

\begin{center}
    \frame{\includegraphics[width=0.45\textwidth]{lecture-07/time-scaling-DT}}
\end{center}

Because we're working with DT signals, it's hard to formally put into words
the time-scaling property.

Other than time-scaling and duality, all of the other CT Fourier transform properties
can be ported over to DT easily.
(Proofs are obvious.)

\begin{theorem}{Even / Odd}
    If $f[n]$ is real and even, $F(\Omega)$ is real and even.
    If $f[n]$ is real and odd, $F(\Omega)$ is imaginary and odd.
\end{theorem}

\begin{theorem}{Hermitian Symmetry}
    If $f[n]$ is real,
    then $F(-\Omega) = \overline{F(\Omega)}$ for all $\Omega$.
\end{theorem}

\begin{theorem}{DT Areas}
    The value of $f[0]$ is $\frac{1}{2\pi} \int_0^{2\pi} F(\Omega) \, \mathrm{d}\Omega$,
    and the value of $F(0)$ is $\sum_{n = -\infty}^{\infty} f[n]$.
\end{theorem}

\begin{theorem}{Time Delay}
    The FT of $f[n - n_0]$ is $e^{-j \Omega n_0} F(\Omega)$.
\end{theorem}

\begin{theorem}{Derivative FT}
    If the FT of $f[n]$ is $F(\Omega)$,
    then the FT of $g[n] := n f[n]$ is $G(\Omega) = j F'(\Omega)$.
\end{theorem}

\begin{theorem}{Kronecker Delta FT}
    Let the \textbf{Kronecker delta function} $\delta[n]$
    be defined so that $\delta[0] = 1$ and
    $\delta[k] = 0$ for all $k \neq 0$.
    The FT of $\delta[n]$ is the constant function $F(\Omega) = 1$.
\end{theorem}

\begin{theorem}{Periodic DT}
    Suppose $f[n]$ is a periodic signal with Fourier coefficients $F[k]$,
    so that $f[n] = \sum_{k = 0}^{N - 1} F[k] e^{jn(2\pi k/N)}$.
    Then, treating $F[k]$ as an infinite array with period $N$,
    the Fourier transform of $f[n]$ is
    $F(\Omega) = \sum_{k = -\infty}^{\infty} 2\pi F[k] \delta\left(\Omega - \frac{2\pi k}{N}\right)$.
\end{theorem}

Here's what that last theorem looks like visually.

\begin{center}
    \frame{\includegraphics[width=0.3\textwidth]{lecture-07/delta-DT}}
\end{center}

\end{document}